%% kapa: exam=1 round=32 year=2021 period=1 form=A subject=civil_law subject_ko=민법 seq=16 number=16
%% source: Q-Net 감정평가사 기출문제 (cst00309) artlSeq=5208397 file=extracted/1008/5208397/A형/2021년도 제32회 감정평가사 1차 1교시 A형.pdf
%% conversion: kordoc 4.24.51 -> tools/gen1.py (auto); review=proofread-sonnet; flags=none
\begin{question}{16}
표현대리에 관한 설명으로 옳지 않은 것을 모두 고른 것은? (다툼이 있으면 판례에 따름)
\begin{framedtext}
ㄱ. 대리권 소멸 후의 표현대리에 관한 규정은 임의대리에만 적용된다. \par
ㄴ. 표현대리를 주장할 때에는 무권대리인과 표현대리에 해당하는 무권대리 행위를 특정하여 주장하여야 한다. \par
ㄷ. 강행법규를 위반하여 무효인 법률행위라 하더라도 표현대리의 법리는 준용될 수 있다. \par
ㄹ. 표현대리가 성립하는 경우에도 상대방에게 과실이 있다면 과실상계의 법리를 유추적용하여 본인의 책임을 경감할 수 있다. \par
\end{framedtext}
\begin{choices}
  \item ㄱ, ㄴ
  \item ㄴ, ㄷ
  \item ㄱ, ㄴ, ㄷ
  \item ㄱ, ㄷ, ㄹ
  \item ㄴ, ㄷ, ㄹ
\end{choices}
\end{question}
\answer{4}
